
\documentclass{article}
\usepackage{colortbl}
\usepackage{makecell}
\usepackage{multirow}
\usepackage{supertabular}

\begin{document}

\newcounter{utterance}

\centering \large Interaction Transcript for game `cladder', experiment `full\_v1.5\_default', episode 2060 with llama{-}3.1{-}8b{-}instruct.
\vspace{24pt}

{ \footnotesize  \setcounter{utterance}{1}
\setlength{\tabcolsep}{0pt}
\begin{supertabular}{c@{$\;$}|p{.15\linewidth}@{}p{.15\linewidth}p{.15\linewidth}p{.15\linewidth}p{.15\linewidth}p{.15\linewidth}}
    \# & \multicolumn{2}{c}{Player} && \multicolumn{2}{c}{Game Master} \\
    \hline

    \theutterance \stepcounter{utterance}  
    & & & \multicolumn{4}{p{0.6\linewidth}}{
        \cellcolor[rgb]{0.9,0.9,0.9}{
            \makecell[{{p{\linewidth}}}]{
                \texttt{\tiny{[P1$\langle$GM]}}
                \texttt{You are an expert in causal inference. The following question is not a typical commonsense query, but rather a meticulously designed question created by a professor specializing in causal inference, intended to assess the students' mastery of the course content.} \\
\texttt{Imagine a self{-}contained, hypothetical world with only the following conditions, and without any unmentioned factors or causal relationships: Encouragement level has a direct effect on studying habit and exam score. Studying habit has a direct effect on exam score. The overall probability of encouragement is 91\%. For students who are not encouraged, the probability of high exam score is 72\%. For students who are encouraged, the probability of high exam score is 52\%. Is high exam score more likely than low exam score overall?} \\
\texttt{Start your answer with 'Yes' or 'No', followed by additional reasoning or evidence to support your explanation.} \\
            }
        }
    }
     \\ \\

    \theutterance \stepcounter{utterance}  
    & \multicolumn{4}{p{0.6\linewidth}}{
        \cellcolor[rgb]{0.561,0.737,0.561}{
            \makecell[{{p{\linewidth}}}]{
                \texttt{\tiny{[P1$\rangle$GM]}}
                \texttt{No.} \\
\\ 
\texttt{To answer this question, we need to apply the principles of causal inference, specifically the concept of counterfactuals and potential outcomes. Let's break down the given information:} \\
\\ 
\texttt{1. Encouragement level has a direct effect on studying habit and exam score.} \\
\texttt{2. Studying habit has a direct effect on exam score.} \\
\texttt{3. The overall probability of encouragement is 91\%.} \\
\texttt{4. For students who are not encouraged, the probability of high exam score is 72\%.} \\
\texttt{5. For students who are encouraged, the probability of high exam score is 52\%.} \\
\\ 
\texttt{We are asked to determine whether high exam score is more likely than low exam score overall. To do this, we need to consider the potential outcomes under different scenarios.} \\
\\ 
\texttt{Let's define two potential outcomes:} \\
\\ 
\texttt{{-} High exam score (H)} \\
\texttt{{-} Low exam score (L)} \\
\\ 
\texttt{We can use the law of total probability to calculate the overall probability of high exam score. However, we need to consider the counterfactuals, i.e., the outcomes under different scenarios.} \\
\\ 
\texttt{Let's define two counterfactuals:} \\
\\ 
\texttt{{-} P(H|not encouraged) = 0.72 (probability of high exam score given no encouragement)} \\
\texttt{{-} P(H|encouraged) = 0.52 (probability of high exam score given encouragement)} \\
\\ 
\texttt{Since the overall probability of encouragement is 91\%, we can calculate the probability of high exam score under the scenario where a student is encouraged:} \\
\\ 
\texttt{P(H|encouraged) = 0.91 \textbackslash{}* 0.52 + 0.09 \textbackslash{}* 0.72 (since 91\% of students are encouraged, and 9\% are not)} \\
\\ 
\texttt{P(H|encouraged) = 0.4712 + 0.0648} \\
\texttt{P(H|encouraged) = 0.536} \\
\\ 
\texttt{Now, let's calculate the probability of high exam score under the scenario where a student is not encouraged:} \\
\\ 
\texttt{P(H|not encouraged) = 0.09 \textbackslash{}* 0.72} \\
\texttt{P(H|not encouraged) = 0.0648} \\
\\ 
\texttt{To calculate the overall probability of high exam score, we need to consider the probability of encouragement and the counterfactuals:} \\
\\ 
\texttt{P(H) = P(encouraged) \textbackslash{}* P(H|encouraged) + P(not encouraged) \textbackslash{}* P(H|not encouraged)} \\
\texttt{P(H) = 0.91 \textbackslash{}* 0.536 + 0.09 \textbackslash{}* 0.0648} \\
\texttt{P(H) = 0.48776 + 0.005808} \\
\texttt{P(H) = 0.493568} \\
\\ 
\texttt{Now, let's calculate the overall probability of low exam score:} \\
\\ 
\texttt{P(L) = 1 {-} P(H)} \\
\texttt{P(L) = 1 {-} 0.493568} \\
\texttt{P(L) = 0.506432} \\
\\ 
\texttt{Since P(L) > P(H), we can conclude that low exam score is more likely than high exam score overall.} \\
            }
        }
    }
    & & \\ \\

    \theutterance \stepcounter{utterance}  
    & & & \multicolumn{2}{p{0.3\linewidth}}{
        \cellcolor[rgb]{0.9,0.9,0.9}{
            \makecell[{{p{\linewidth}}}]{
                \texttt{\tiny{[GM$|$GM]}}
                \texttt{No} \\
            }
        }
    }
    & & \\ \\

    \theutterance \stepcounter{utterance}  
    & & & \multicolumn{2}{p{0.3\linewidth}}{
        \cellcolor[rgb]{0.9,0.9,0.9}{
            \makecell[{{p{\linewidth}}}]{
                \texttt{\tiny{[GM$|$GM]}}
                \texttt{yes} \\
            }
        }
    }
    & & \\ \\

    \theutterance \stepcounter{utterance}  
    & & & \multicolumn{2}{p{0.3\linewidth}}{
        \cellcolor[rgb]{0.9,0.9,0.9}{
            \makecell[{{p{\linewidth}}}]{
                \texttt{\tiny{[GM$|$GM]}}
                \texttt{game\_result = LOSE} \\
            }
        }
    }
    & & \\ \\

\end{supertabular}
}

\end{document}
